\documentclass[10pt,a4paper]{amsart}
\usepackage[margin=19mm]{geometry}
\usepackage{amsmath,amssymb,mathtools}
\usepackage[scr=rsfs,cal=euler]{mathalfa}
\usepackage{xfrac}
\usepackage[unicode,hidelinks,bookmarks=false]{hyperref}
\newcommand{\ie}{i.e.\ }
\newcommand{\cf}{cf.\ }
\newcommand{\resp}{respectively\ }
\newcommand{\Subsection}{\subsection}
\providecommand{\nobreakdash}{\nobreak\hbox{-}}
\setlength{\emergencystretch}{2em}
\multlinegap0pt
\usepackage{amssymb}
\usepackage{tikz}
\usepackage{float}
\usepackage[shortlabels]{enumitem}
\usepackage[all]{xy}
\newtheorem{proposition}{Proposition}
\newtheorem{definition}{Definition}
\newtheorem{lemma}{Lemma}
\DeclareMathOperator{\Lip}{Lip}
\newcommand{\de}{\partial}
\DeclareMathOperator{\diam}{diam}
\DeclareMathOperator{\dist}{dist}
\newcommand{\ve}{\varepsilon}
\title{Interior control for surfaces with positive scalar curvature and its application}
\author{Shuli Chen, Jianchun Chu, Jintian Zhu}
\date{Source: arXiv:2502.14216v3 (CC BY 4.0)}
\begin{document}
\maketitle

\noindent\emph{Source and licence.} Interior control for surfaces with positive scalar curvature and its application. Version arXiv:2502.14216v3 (CC BY 4.0), distributed by arXiv under the Creative Commons Attribution 4.0 International licence, http://creativecommons.org/licenses/by/4.0/, as recorded in the arXiv licence metadata for this exact version and shown on the abstract page https://arxiv.org/abs/2502.14216v3. Full licence notice: \texttt{licenses/arxiv-2502.14216-CC-BY-4.0.txt}.\par\smallskip

\noindent\textit{Editorial scope.} Counted unit: the article's proposition interior (source0.tex lines 154--170). The article's complete original proof of it is delivered with it (source0.tex lines 319--480, one \texttt{proof} environment of 162 lines, which the article places away from the statement). No mathematics is written, repaired or re-derived: the statement and the proof are the article's own lines, in the article's own notation, and no retained line is substituted or re-flowed.  Retained uncounted as the source-local context the proof needs: lines 282--285; lines 292--305; lines 308--314.  Nothing was omitted from any retained span, so the package carries no omission record.  The retained lines cite 2 works, whose entries are transcribed verbatim from the article's own bibliography source in the e-print and delivered with short editorial labels recorded in the item's reference mapping.  No retained line contains a figure environment, an image-inclusion command or drawing code, so no image is delivered and the package declares no asset dependency.  Editorial formatting only, in addition: an AMS \texttt{amsart} preamble that restates the article's own declarations and macros used by the retained lines, namely the environments \texttt{proposition}, \texttt{definition}, \texttt{lemma} on separate counters, together with the standard packages those lines need.  The retained environments print in this document as Proposition 1, Definition 1, Lemma 1 in order of appearance, whereas the article prints the same items with the numbering of its own full document.  The complete native source of the article as distributed in the arXiv e-print for this exact version is bundled unchanged as \texttt{sources/173569/source0.tex}.
\begin{definition}\label{Defn: stabilized scalar curvature}
Let $(M,g)$ be a Riemannian manifold possibly with non-empty boundary and $R$ a smooth function on $M$. We say that $(M,g)$ has $\mathbb T^l$-stabilized scalar curvature $R$ if there are $l$ positive smooth functions $v_1,\ldots,v_l$ on $M$ such that the scalar curvature $\tilde R$ of $(M\times \mathbb T^l,g+\sum_i v_i^2\mathrm d\theta_i^2)$ is the $\mathbb T^l$-invariant extension of $R$.
We say that $(M,g)$ has $\mathbb T^*$-stabilized scalar curvature $R$ if $(M,g)$ has $\mathbb T^l$-stabilized scalar curvature $R$ for some positive integer $l$.
\end{definition}

\begin{lemma}[Lemma 2.12 of \cite{SWWZ24}]\label{Lem: D estimate}
Given any positive constant $c_0$, there is a positive constant $D_{0}=D_0(c_0)$ such that the following holds.
Let $(M^{n},\partial_\pm,g)$ be a Riemannian band with $n\leq7$ and $\overline{\partial M-(\partial_+\cup\partial_-)}=\emptyset$. If
\[
\rho: (M,\de_{\pm}) \to ([-D,D],\pm D)
\]
is a smooth function satisfying
\begin{itemize}\setlength{\itemsep}{1mm}
\item $\Lip\rho<1$;
\item $(M,g)$ has $\mathbb T^*$-stabilized scalar curvature $R$ such that $R\geq c_{0}$ in $\rho^{-1}([-1,1])$ and $R\geq0$ in $M$;
\item Any closed hypersurface $\Sigma$ in $M$ separating $\de_{-}$ from $\de_{+}$ does not admit any metric with positive $\mathbb{T}^{*}$-stabilized scalar curvature,
\end{itemize}
then $D\leq D_{0}$.
\end{lemma}

\begin{lemma}[Lemma 2.14 of \cite{CCZ23}]\label{Lem: free boundary}
Let $(M^{n},\partial_\pm,g)$ be a Riemannian band with $2\leq n\leq 7$ and $\dist(\partial_+,\partial_-)>2d_0$. If $\Gamma:=\overline{\partial M-(\partial_+\cup\partial_-)}$ does not intersect $\de_{\pm}$ and $(M,g)$ has $\mathbb T^l$-stabilized scalar curvature $R$, then there exists a hypersurface $\Sigma$ (possibly with non-empty boundary $\partial\Sigma$) intersecting $\Gamma$ orthogonally such that
\begin{itemize}\setlength{\itemsep}{1mm}
\item $\Sigma$ and $\partial_-$ bound a relative region $\Omega$ relative to $\Gamma$, i.e. $\partial\Omega-(\Sigma\cup \partial_-)\subset \Gamma$;
\item $(\Sigma,g)$ has $\mathbb T^{l+1}$-stabilized scalar curvature $R'$ where $R'\geq R-\pi^2/d_0^2$.
\end{itemize}
\end{lemma}

\begin{proposition}\label{Prop: interior}
     Given a positive and continuous function $L:[0,+\infty)\to (0,+\infty)$ there is a positive function $K:(0,+\infty)\to (0,+\infty)$ such that the following property holds: for any $s_{0}>0$, if $(\Sigma,g,\rho)$ is a triple  satisfying
    \begin{itemize}\setlength{\itemsep}{1mm}
        \item $(\Sigma,g)$ is a connected Riemannian surface with non-empty boundary and $\rho:\Sigma\to[0,+\infty)$ is a smooth and proper function with $\Lip\rho <1$;
        \item

        $\rho>K(s_0)$ on $\partial\Sigma$;



        \item  $(\Sigma,g)$ has $\mathbb{T}^*$-stabilized scalar curvature $R\geq L\circ \rho$ in $\Sigma_{K(s_0)}$, where $\mathbb{T}^*$-stabilized scalar curvature is defined in Definition \ref{Defn: stabilized scalar curvature} and we use the notation $\Sigma_s:=\rho^{-1}([0,s])$,
    \end{itemize}
    then we can find finitely many pairwise disjoint disks $\{D_j\}$ in $\Sigma_{K(s_0)}$ such that
    $$\Sigma_{s_0}\subset\bigcup_j D_j$$
    and for each connected component $\Sigma_{s_0}^{o}$ of $\Sigma_{s_0}$, the extrinsic diameter of $\Sigma_{s_0}^{o}$ in $(\Sigma_{K(s_0)},g)$ is bounded by
   $$ 4\pi \left(\min_{[0,K(s_0)]}L\right)^{-\frac{1}{2}}.$$
\end{proposition}

\begin{proof}[Proof of Proposition \ref{Prop: interior}]
It suffices to determine the value of $K(s_{0})$ for $s_{0}$. Set
\[
c_{0} = \min_{[0,s_{0}+2]}L,\ \
s_{1} = s_{0}+D_{0}(c_{0})+D_{0}(3c_{0}/4)+2,
\]
\[
c_1=\min_{[0,s_{1}+3]}L,\ \ K(s_{0})=s_1+D_0(c_1)+3,
\]
where $D_0(c_0)$, $D_0(3c_0/4)$ and $D_0(c_1)$ are the constants defined in Lemma \ref{Lem: D estimate}.
Let $s'$ be a regular value of $\rho$ in $(K(s_0)-1,K(s_0))$.
We deal with interior topology and interior geometry separately.


\bigskip

\noindent
{\bf I. Interior topology.}
The argument is broken into three steps.

\medskip

\noindent
{\bf Step 1.} Let $s_{1}^{*}$ be a regular value of $\rho$ in $(s_{1},s_{1}+1)$. Then each component $\Sigma_{s_1^*}^o$ of $\Sigma_{s_1^*}$ is homeomorphic to a $2$-sphere with finitely many disks removed.

\medskip

Otherwise, some component $\Sigma_{s_1^*}^o$ must be $m\mathbb{T}^{2}$ or $m\mathbb{RP}^{2}$ after filling with disks. In particular, $\Sigma_{s_1^*}^o$ has the form of $\mathbb{T}^2\#S$ or $\mathbb{RP}^2\#S$ for some surface $S$, then we can find $\hat \Sigma^o_{s_1^*}\subset \Sigma_{s_1^*}^o$ homeomorphic to $\mathbb T^2-B$ or $\mathbb{RP}^{2}-B$.
Let $\Sigma_{s'}^o$ denote the component of $\Sigma_{s'}$ containing $\Sigma_{s_1^*}^o$, and let $S_+$ and $S_-$ denote copies of $\Sigma_{s'}^o\setminus \hat\Sigma_{s_1^*}^o$. Consider a double cover $\tilde\Sigma_{s_1^*}^o$ of $\hat\Sigma_{s_1^*}^o$ and clearly $\tilde\Sigma_{s_1^*}^o$ has the form of $\mathbb T^2-(B_+\sqcup B_-)$ or $\mathbb S^2-(B_+\sqcup B_-)$, where $B_+$ and $B_-$ are lifts of $B$. Correspondingly, the surface $\Sigma_{s'}^o$ is lifted to some double cover
\[
\tilde\Sigma_{s'}^o=S_+\cup \tilde\Sigma_{s_1^*}^o \cup S_-.
\]
Denote $\pi:\tilde\Sigma_{s'}^o\to \Sigma_{s'}^o$ to be the covering map. Then the map
\begin{equation*}
\tilde \rho(x)=\left\{
\begin{array}{cc}
(\rho\circ \pi-s_1^*)_+,& x\in S_+;\\[1mm]
-(\rho\circ \pi-s_1^*)_+,& x\in S_-;\\[1mm]
0,&\mbox{otherwise},
\end{array}\right.
\end{equation*}
is a Lipschitz function on $\tilde\Sigma_{s'}^o$ with $\Lip \tilde \rho<1$. Also we see that the function $\tilde \rho $ has range $[s_1^*-s',s'-s_1^*]$, and that $\tilde \Sigma_{s'}^o$ has $\mathbb{T}^*$-stabilized scalar curvature $R\geq c_1\chi_{[-2,2]}\circ\tilde \rho$. We further notice that the third hypothesis of Lemma \ref{Lem: D estimate} is trivial on surfaces since $\mathbb{S}^1$ does not admit any metric with positive $\mathbb{T}^{*}$-stabilized scalar curvature. Thus by smooth approximation and Lemma \ref{Lem: D estimate} we obtain
$$
D_0(c_1)+1 = K(s_{0})-1-(s_{1}+1) < s'-s_1^*\leq D_0(c_1),
$$
which leads to a contradiction.


\medskip

\noindent
{\bf Step 2.} Let $s_0^*$ be a regular value of $\rho$ in $(s_0,s_0+1)$. Then each circle $\gamma$ in $\rho^{-1}(s_0^*)$ bounds a disk in $\Sigma_{s_{1}^{*}}$.

\medskip

Let $\Sigma_{s_1^*}^o$ be the component of $\Sigma_{s_1^*}$ containing $\gamma$. Suppose that $\gamma$ does not bound a disk in $\Sigma_{s_1^*}$. Recall that $\Sigma_{s_1^*}^{o}$ is homeomorphic to a $2$-sphere with finitely disks removed. From this we conclude that $\gamma$ separates the boundary components of $\partial\Sigma_{s_1^*}^o$ into two non-empty collections. Denote
$$\Sigma_{s_1^*}^o-\gamma=\Omega_+\sqcup \Omega_-.$$
Define
\begin{equation*}
\rho_{s_1^*}(x)=\left\{
\begin{array}{cc}
\min\{\dist(x,\gamma), s_1^*-s_0^*\},& x\in \Omega_+;\\[1mm]
-\min\{\dist(x,\gamma), s_1^*-s_0^*\},& x\in \Omega_-;\\[1mm]
0,&\mbox{otherwise}.
\end{array}\right.
\end{equation*}
Then $\rho_{s_1^*}$ is a Lipschitz function on $\Sigma_{s_1^*}^o$ with $\Lip \rho_{s_1^*}\leq1$. Also $\rho_{s_1^*}$ has range $[s_0^*-s_1^*,s_1^*-s_0^*]$, and that $\Sigma_{s_1^*}^o$ has $\mathbb{T}^*$-stabilized scalar curvature $$R\geq c_0\chi_{[-s_0-2+s_0^*,s_0+2-s_0^*]}\circ \rho_{s_1^*}.$$
By smooth approximation and Lemma \ref{Lem: D estimate} we obtain
$$
D_0(c_0)+D_0(3c_0/4)+1 = s_{1}-s_{0}-1 < s_1^*-s_0^*\leq D_0(c_0),
$$
which leads to a contradiction.

\medskip

\noindent
{\bf Step 3.} $\Sigma_{s_0}$ is contained in $\sqcup_j D_j$, where $\{D_j\}$ are finitely many disjoint disks in $\Sigma_{s_1^*}$.

\medskip

By passing to each component we can assume $\Sigma_{s_1^*}$ to be connected. Due to the fact $\Sigma_{s_0}\subset \Sigma_{s_0^*}$ it suffices to show that $\Sigma_{s_0^*}$ is contained in $\sqcup_j D_j$, where $\{D_j\}$ are finitely many disjoint disks in $\Sigma_{s_1^*}$. From our assumption we know that $\Sigma$ is connected and so $\Sigma_{s'}$ has non-empty boundary. In particular, $\Sigma_{s_1^*}$ has non-empty boundary as well. Denote $\rho^{-1}(s_0^*)=\sqcup_k\gamma_k$. From the previous steps we conclude that each $\gamma_k$ bounds a unique disk $D_k$ in $\Sigma_{s_1^*}$.

We claim that for $k\neq l$ one of the following happens:
\begin{itemize}\setlength{\itemsep}{1mm}
\item $D_k$ and $D_l$ are disjoint;
\item $D_k\subset D_l$ or $D_l\subset D_k$.
\end{itemize}
Otherwise, we can find a point $p_{kl}\in D_k\cap D_l$ and a pair of points $p_k\in D_k\setminus D_l$ and $p_l\in D_l\setminus D_k$. Connecting $p_{kl}$ and $p_k$ with path in $D_k$ we see $\gamma_l\cap D_k\neq\emptyset$. Since $\gamma_k$ and $\gamma_l$ are disjoint, we conclude $\gamma_l\subset D_k$ and similarly $\gamma_k\subset D_l$. In particular, the union $D_k\cup D_l$ is a closed surface in $\Sigma_{s_1^*}$, which is impossible since $\Sigma_{s_1^*}$ is connected with non-empty boundary.

Now we can take the collection of maximal disks $\{D_{j,\mathrm{max}}\}$ with respect to the inclusion. From previous claim we conclude that $D_{j,\mathrm{max}}$ are pairwise disjoint.

It remains to show $\Sigma_{s_0^*}\subset \sqcup D_{j,\mathrm{max}}$. Take any component $\Sigma_{s_0^*}^o$ of $\Sigma_{s_0^*}$. Since $\Sigma_{s_0^*}^o$ does not cross $\sqcup_k\gamma_k$, either $\Sigma_{s_0^*}^o\subset \sqcup D_{j,\mathrm{max}}$ or $\Sigma_{s_0^*}^o$ is disjoint with all disks $D_k$. Denote $\partial\Sigma_{s_0^*}^o=\sqcup_l \gamma_{k_l}$. Then the union
$$
\Sigma_{s_0^*}^o\cup\left(\bigsqcup_lD_{k_l}\right)
$$
is a closed surface in $\Sigma_{s_1^*}$, which leads to the same contradiction as before.


\bigskip
\noindent
{\bf II. Interior geometry.}
To establish the extrinsic diameter estimate, we argue by contradiction. Set
\[
c_2=\min_{[0,K(s_{0})]}L.
\]
If
\[
\diam(\Sigma_{s_0}^{o} \subset (\Sigma_{s_1^*},g)) > 4\pi \left(\min_{[0,K(s_{0})]}L\right)^{-\frac{1}{2}} = \frac{4\pi}{\sqrt{c_{2}}},
\]
then there exist two points $p_{+},p_{-}\in\Sigma_{s_0}^{o}$ such that $\dist(p_{+},p_{-})>4\pi/\sqrt{c_{2}}$. Recall that $s'$ is a regular value of $\rho$ in $(K(s_{0})-1,K(s_{0}))$. For sufficiently small $\ve>0$, set
\[
M = \Sigma_{s'}\setminus(B_{\ve}^{g}(p_{-})\cup B_{\ve}^{g}(p_{+})), \ \
\de_{-} = \de B_{\ve}^{g}(p_{-}), \ \
\de_{+} = \de B_{\ve}^{g}(p_{+}), \ \
\Gamma = \de\Sigma_{s'},
\]
and then $(M,\de_{\pm},g)$ is a Riemannian band with $\dist(\partial_+,\partial_-)>4\pi/\sqrt{c_{2}}$ and $\Gamma\cap(\de_{+}\cup\de_{-})=\emptyset$. Recall that $(\Sigma,g)$ has $\mathbb{T}^l$-stabilized scalar curvature $R\geq L\circ \rho$ in $\Sigma_{K(s_0)}$ and so in $\Sigma_{s'}$.
By Lemma \ref{Lem: free boundary} (with the choice $d_{0}=2\pi/\sqrt{c_{2}}$), there exists a curve $\gamma$ satisfying
\begin{itemize}\setlength{\itemsep}{1mm}
\item[(a)] $\gamma$ and $\partial_-$ bound a relative region $\Omega$ relative to $\Gamma$;
\item[(b)] $\gamma$ has the positive $\mathbb T^{l+1}$-stabilized scalar curvature
\[
R' \geq R-\frac{c_{2}}{4} \geq L\circ \rho-\frac{c_{2}}{4} \geq c_{2}-\frac{c_{2}}{4} = \frac{3}{4}c_{2} > 0.
\]
\end{itemize}
We split the argument into two cases:

\medskip

{\it Case 1. $\gamma\cap\rho^{-1}(s')=\emptyset$.} In this case, $\de\gamma=\emptyset$ and each component of $\gamma$ is $\mathbb{S}^{1}$. Then (b) shows that $\mathbb{S}^{1}$ has positive $\mathbb T^{l+1}$-stabilized scalar curvature. This contradicts the fact that $\mathbb{T}^{n}$ cannot admit any smooth metric with positive scalar curvature.

\medskip

{\it Case 2. $\gamma\cap\rho^{-1}(s')\neq\emptyset$.} Let $\sigma$ be a path in $\Sigma_{s_{0}}^{o}$ such that
\begin{itemize}\setlength{\itemsep}{1mm}
\item $\sigma$ connects $p_{-}$ and $p_{+}$;
\item $\sigma$ has non-zero algebraic intersection with $\de_{-}$.
\end{itemize}
Using (a), $\gamma$ is relative homologous to $\de_{-}$ relative to $\Gamma$ and then $\gamma$ has non-zero algebraic intersection with $\sigma$. In particular, we obtain $\gamma\cap\sigma\neq\emptyset$ and so $\gamma\cap\Sigma_{s_{0}}\neq\emptyset$. Fix a point $x_{0}\in\gamma\cap\Sigma_{s_{0}}$ and the component of $\gamma$ containing $x_{0}$. For convenience, we still denote this component by $\gamma$. Choose an arc length parametrization of $\gamma$ such that
\[
\gamma:[-D',D'']\to M, \ \
\gamma(0) = x_{0}, \ \
\gamma(-D'),\,\gamma(D'') \in \rho^{-1}(s'),
\]
where $D'$ and $D''$ are two positive constants. We next assume that $D'\geq D''$, as the case $D''>D'$ can be proved similarly. Since $\Lip\rho <1$, then
\begin{equation}\label{rho gamma t control}
|\rho(\gamma(t))-\rho(\gamma(0))| < |t| \ \text{for $t\in[-D',D''] ,\, t \neq 0$}.
\end{equation}
Choosing $t=D''$ and recalling the facts $\gamma(D'') \in \rho^{-1}(s')$ and $\rho(\gamma(0))\leq s_0$, we can derive from \eqref{rho gamma t control} that
\begin{equation}\label{D'' lower bound}
D'' > s' - s_{0} > K(s_{0})-2-s_{0} = D_{0}(c_{0})+D_{0}(3c_{0}/4)+D_0(c_1)+3.
\end{equation}
Using \eqref{rho gamma t control} again, we obtain
\[
\rho(\gamma(t)) < s_{0}+1 \ \text{for $t\in[-1,1]$},
\]
which implies $\gamma([-1,1])\subset\Sigma_{s_{0}+1}$. Together with (b) and $c_{2}\leq c_{0}$, we know that $\gamma$ has positive $\mathbb T^{l+1}$-stabilized scalar curvature $R'$ with
\[
R' \geq R-\frac{c_{2}}{4} \geq L\circ \rho-\frac{c_{2}}{4} \geq c_{0}-\frac{c_{2}}{4} \geq \frac{3c_{0}}{4} > 0 \ \text{in $\gamma([-1,1])$}.
\]
Applying Lemma \ref{Lem: D estimate} to $\gamma([-D'',D''])$ and $\rho|_{\gamma([-D'',D''])}$, we obtain $D''\leq D_{0}(3c_{0}/4)$, which contradicts with \eqref{D'' lower bound}.
\end{proof}

\begin{thebibliography}{99}
\bibitem[CCZ23]{CCZ23} Chen, S., Chu, J., Zhu, J. {\em Positive scalar curvature metrics and aspherical summands}, preprint, arXiv: 2312.04698.
\bibitem[SWWZ24]{SWWZ24} Shi, Y., Wang, J., Wu, R., Zhu, J. {\em On open manifolds admitting no complete metric with positive scalar curvature}, preprint, arXiv: 2404.01660.
\end{thebibliography}
\end{document}
